\begin{figure}[H]\centering
\begin{tikzpicture}[x=1mm,y=1mm]
\node[box,text width=39mm] at (28,0) {Launch client};
\node[box,text width=39mm] at (130,0) {Away service};
\draw[dashed,gray] (28,-7)--(28,-83);
\draw[dashed,gray] (130,-7)--(130,-83);
\draw[flow] (28,-15)--node[above,font=\footnotesize]{GET info; verify identity}(130,-15);
\draw[flow] (28,-26)--node[above,font=\footnotesize]{WebSocket upgrade}(130,-26);
\draw[flow] (130,-37)--node[above,font=\footnotesize]{hello: boot + retained range}(28,-37);
\draw[flow] (28,-48)--node[above,font=\footnotesize]{subscribe: mission + resume cursor}(130,-48);
\draw[flow] (130,-59)--node[above,font=\footnotesize]{ready; replay OR reset + snapshot}(28,-59);
\draw[flow] (130,-70)--node[above,font=\footnotesize]{sync.complete; ordered live events}(28,-70);
\draw[flow,<->] (28,-81)--node[above,font=\footnotesize]{heartbeat / acknowledgement}(130,-81);
\end{tikzpicture}
\caption{Synchronization has an atomic event barrier. A snapshot repairs
current state; a reset explicitly records unavailable history.}
\end{figure}
